% !TEX root = handbook-physics.tex
%!TEX TS-program = pdflatex
%!TEX encoding = UTF-8 Unicode






\chapter{Dielectrics In An Electric Field}
\label{ch-18}

\section{Dipole Moments of Molecules of a Dielectric}
\label{sec-18.1}

\subsection{}\label{18.1.1} \redem{Dielectrics} are substances that do not conduct electric current. At not especially high temperatures and under conditions when dielectrics are not subject to very strong electric fields, these substances, in contrast to conductors, have no carriers of electric current.

\subsection{}\label{18.1.2} The molecules of a dielectric are electrically neutral and contain an equal number of positive and negative charges. Nevertheless, the molecules do have electrical properties. As a first approximation, a molecule of a dielectric can be regarded as a dipole having the electric dipole moment $\vb{p}_{e} =q \vb{l}$ \ssect{15.2.3}, where $q$ is the magnitude of the total positive (and also total negative) charges, located, respectively, at the centres of gravity of the charges, $\vb{l}$ is the distance between the centres of gravity of the positive and negative charges. As any dipole, a molecule of matter sets up an electric held in the surrounding space \ssect{15.2.3}.

\subsection{}\label{18.1.3} A dielectric is said to be \redem{nonpolar} (a \redem{nonpolar dielectric}) if the electrons of the atoms in its molecules are arranged symmetrically with respect to their nuclei (\ce{H2}, \ce{O2}, \ce{CCl4}, etc.). In such molecules the centres of gravity of the positive and negative charges coincide when there is no external electric held [$l = 0$ \ssect{18.1.2}] and the dipole moment $\vb{p}_{e}$ of the molecule is equal to zero. When a nonpolar dielectric is placed into an external electric held, the electron shells \ssect{39.3.6} in the atoms (or molecules) are deformed and the centres of gravity of the positive and negative charges are displaced with respect to each other ($l \neq  0$). At this an \redem{induced electric dipole moment} appears in atoms (or molecules) of the dielectric. This moment is proportional to the strength $\vb{E}$ of the external electric held:
\begin{align*}
\vb{p}_{e} & = \varepsilon_{0} \alpha \vb{E} && \text{(in SI units.)} \\
\vb{p}_{e} & =  \alpha \vb{E}&& \text{(in cgse units.)}
\end{align*}
where $\alpha$ is the \redem{polarizability} of the molecule (or atom) and $\varepsilon_{0}$ is the electric constant in SI units \ssect{14.2.7}.


The polarizability of a molecule depends only on its volume. It is important that $\alpha$ does not depend upon temperature. The thermal motion of the molecules in non-polar dielectrics has no effect on the induction of dipole moments in them. Molecules with such dipole moments are like \redem{quasi-elastic (induced)} dipoles.

\subsection{}\label{18.1.4} A \redem{polar dielectric} is one in which the molecules (or atoms) have electrons arranged asymmetrically with respect to their nuclei (\ce{H2O}, \ce{HCl}, \ce{NH3}, \ce{CH3Cl}, etc.). In such molecules the centres of gravity of the positive and negative charges do not coincide, but are at a practically constant distance I from each other. With respect to electrical properties, molecules of polar dielectrics are like \redem{rigid dipoles} that have a \redem{permanent dipole moment} $\vb{p}_{e} = \text{const}$.

\subsection{}\label{18.1.5} A \redem{permanent dipole} placed in a homogeneous external electrostatic field is subject to the action of a force couple \ssect{4.1.6} with the moment of force (torque) $\vb{M}$ equal to $\vb{M} = \cpr{\vb{p}_{e}}{\vb{E}}$.

The moment $\vb{M}$ of the force couple is perpendicular to the plane passing through vectors pe and E. From the tip of vector M rotation through the smaller angle from vectors pe to E is seen to he counterclockwise. In Fig. 18.1 the moment of force M
is directed normal to the diagram into the page and tends to turn the dipole so that pe and E are parallel to each other. In real molecules of polar dielectrics, the molecules are deformed in addition to the turning of the dipole axes into alignment with the field. This leads to the induction of a certain dipole moment \ssect{18.1.3}.

\subsection{}\label{18.1.6} When a permanent dipole is placed into a nonhomo- geneous electric field, in which the field strength E varies along the length Z of the dipole, then, in addition to the torque M \ssect{18.1.5}, the dipole is subject to force F, which is equal to

or, in the more general form:
r= dE dE dE Pex dx ' Pey dy ~T~Pez dz 1
251
where pex,peiJ and pez are the projections of vector pe on the axes of a Cartesian coordinate system. In particular, if vector E is directed along the OX axis and depends only on the coor­ dinate x [Ex = E (x), Ey = Ez = 0 and E = Exi], then
I3y the action of force F a free permanent dipole tends to move to a region of highest field strength.

\subsection{}\label{18.1.7} The potential energy Ev of a permanent dipole in an external electric field is
Ep = —peE = —peE cos 0,
where pe is the electric dipole moment, E is the field strength at the point where the dipole is located and G is the angle be­ tween the dipole axis and vector E. The minus sign indicates that the stable portion of the dipole corresponds to the angle 0 = 0, at which the potential energy of the dipole has its mini­ mum value and vector pe is directed along vector E.

\section{Polarization of Dielectrics}
\label{sec-18.2}

\subsection{}\label{18.2.1} When a polar dielectric \ssect{17.1.4} is not in an external electric field, the vectors of its dipole moments are chaotically oriented as a result of the random thermal motion of the mole­ cules. Consequently, the sum of the dipole moments of all the molecules is equal to zero in any physically infinitesimal volume AT*.
No dipole moments of the molecules whatsoever can be induced \ssect{18.1.3} in a nonpolar dielectric that is not in an external electric field.
* Herev0, where v0 is the volume of a single molecule,
and an immense number of molecules are contained in vol­ ume A7.

\subsection{}\label{18.2.2} When a dielectric is placed in an external electric field, it is polarized. This means that in any element of volume A7, the total dipole moment of its molecules is no longer equal to zero. A dielectric in such a state is said to be polarized (polar­ ized dielectric). In accordance with the structure of the mole­ cules (or atoms), distinction is made between the following three types of polarization:
\begin{enumerate}[label=(\alph*)]
\item Orientation polarization of polar dielectrics \ssect{18.1.4}. The external electric field tends to align the dipole moments of the permanent dipoles along the direction of the field \ssect{18.1.5}. This is hindered by the random thermal motion of the molecules which tends to disarrange the dipoles arbitrarily. The combined action of the field and thermal motion results in preferred orien­ tation of the electric dipole moments along the field. This polarization increases with the strength of the electric field and with a drop in temperature.
\item Electronic polarization in nonpolar dielectrics \ssect{18.1.3} In the molecules of dielectrics of this type dipole moments directed along the field are induced \ssect{18.1.3} by the action of the external electric field. The thermal motion of the molecules has no effect on electronic polarization. In gaseous and liquid dielec- tries, electronic polarization takes place practically simulta­ neously with orientational polarization.
\item Ionic polarization in solid dielectrics having ionic crystal lattices \ssect{40.1.3}, such as NaCl and CsCl. The external electric field causes displacement of all the positive ions in the direction of the field strength E and all the negative ions in the opposite direction.
\end{enumerate}

\subsection{}\label{18.2.3} A quantitative measure of the polarization of a dielec­ tric is the polarization vector Pe. The polarization vector (or simply the polarization) is the ratio of the electric dipole moment of a small volume AV of the dielectric to this volume:
n
where pei- is the electric dipole moment of the ith molecule, and n is the total number of molecules in the volume AV. This volume should be so small that the electric field \ssect{15.1.2} within it can be assumed uniform. At the same time, the number n

of molecules in it should be sufficiently large to allow statistical methods of investigation \ssect{8.2.2} to be applied.

\subsection{}\label{18.2.4} For a homogeneous nonpolar dielectric \ssect{18.1.3} in a uni­ form electric field
Pe= noPe>
where nQis the number of molecules in unit volume, and pe is the dipole moment of a single molecule. Making use of the
formula for pe \ssect{18.1.3} we obtain
Pe = /20e0aE = e0eE (in SI units), Pe = n0aE = eE (in cgse units),
where e = n0a is the dielectric (or electric) susceptibility of the substance.

\subsection{}\label{18.2.5} For a homogeneous polar dielectric \ssect{18.1.4} in a uniform electric field, Pe = n0 (pe), where (pe) is the average value of the component along the field of the permanent dipole moment of the molecule. When a polar dielectric is in a weak external electric fjeld, the dielectric susceptibility is calculated by the Debye-Langevin formula:
n0P%
**=-kr (;nSIunits)-
noPe ... Xe“1W (mcgseunlts)-
Here k is Boltzmann’s constant \ssect{8.4.5}, and T is the absolute temperature. The remaining notation is given in \ssect{18.1.3}. Shown in Fig. 18.2 is the dependence e = Xe (1IT) for non­ polar (a) and for polar (b) molecules. The straight line (b) does not pass through the origin of coordinates because orientational and electronic polarization \ssect{18.2.2b} usually occur in polar molecules and the dielectric susceptibility consists of two parts: e = Xq + Xe’where e an( Xe are expressed by the equations in Sections \ssect{18.2.4} and \ssect{18.2.5}.

\subsection{}\label{18.2.6} If a dielectric is in a uniform electric field, the electrical neutrality of any volume AV of the dielectric, containing a suf­ ficiently large number of molecules, is provided for by the mutual

compensation of the dipole charges of opposite sign alongside one another (Fig. 18.3). In the thin layers at surfaces and *S2 of the dielectric, bounding its volume, surface-bound polarization charges appear as a result of the polarization of the dielectric. At the surface where the lines of force \ssect{15.1.6} of the external electric field enter the dielectric, negative charges of the ends
Fig. 18.2 Fig. 18.3
of the dipole molecules appear. Positive charges appear at the opposite surface. The surface density \ssect{15.2.3} ap of the polar­ ization charges is calculated by the formula
aP= en>
where Pen is the projection of vector Pe on the outward normal to the surface of the dielectric.
In a nonuniform electric field, the polarization of a dielectric is also nonuniform and the polarization vector Pe depends upon the coordinates. Hence, in addition to surface-bound charges, space polarization charges also appear. These are distributed with the volume charge density \ssect{15.2.3} pp:
pp = —div Pe,
where div Pe-- -■ey  ez is the divergence of dx dy dz
the polarization vector.

\section{Relation Between Displacement, Field Strength and Polarization Vectors}
\label{sec-18.3}

\subsection{}\label{18.3.1} Distinction is made in substances between two kinds of charges: free and bound.
Bound charges are those included in the composition of atoms and molecules, as well as the charges of ions in crystalline solids with an ionic lattice \ssect{40.1.3}. Charges not bound to the aforementioned particles of matter are said to be free. Free charges include: (a) current-conducting charges in conducting media [conduction electrons in metals and semiconductors \ssect{16.4.1}, holes in semiconductors \ssect{41.10.3}, ions in electrolytes and gases],
(b) surplus charges that are transmitted to a body by various means and violate its electrical neutrality. An example are charges imparted from outside onto the surface of a dielectric. 

\subsection{}\label{18.3.2} Generally, an electric held is set up in a dielectric by both free and bound charges. The held strength vector E characterizes the resultant held in the dielectric, set up by the two kinds of charges, and depends upon its electrical properties: the relative permittivity (dielectric constant) er \ssect{14.2.4}. But the primary source'of an electric held in a dielectric is the free charges. As a matter of fact, a held of bound charges is set up in a dielectric as a result of its polarization when it is placed in an external electric held set up by free electric charges. 

\subsection{}\label{18.3.3} The information contained in Sections \ssect{18.3.1} and \ssect{18.3.2} call for a certain revision of the Ostrogradsky-Gauss theorem
ly
hand side of the theorem is to be regarded as the algebraic
sum of the free electric charges enclosed by surface S, i.e. k
yi qi = <7free»and the theorem is then of the form i=l
DndS = <7free (in SI units),
DdS= DndS=4jtgfree(incgseunits). S

\ssect{15.3.3}. The algebraic sum of the charges V, qt on the right­

In this form the Ostrogradsky-Ganss theorem is valid for an electric field either in homogeneous and isotropic or in nonho- mogeneous and anisotropic media (cf. \ssect{15.3.3}.

\subsection{}\label{18.3.4} For an electric field in vacuum
D = e0E (in SI units),
D = E (in cgse units).
The flux of the field strength vector E through an arbitrary closed surface S in vacuum is
s0EndS= gfree (in SI units), s
EndS= 4itgfree (in cgse units). s
In accordance with Section \ssect{18.3.2}, the flux of vector E for a field in a substance is
§e0EndS = <7free+ gbnd (in SI units), S
 £nd£= 4n(gfree+ 0bnd) (in cgse units). s
The sum of bound charges gbnd within closed surface S is cal­ culated by the equation
s
where Pen is the projection of the polarization vector on the outward normal to surface dS. Hence
§(e„ £n+Pen) = 9frec (in SI units),
 IS.',I ferroelectric materials 257
A correlation with the general formulation of the Ostrogradsky- (iauss theorem \ssect{18.3.3} enables a relation to be established be* Iween vectors D, E and Pe:
oEn+ Pen (in SI units), />„ En + 4jtPen (in cgse units),
or, by virtue of the arbitrary nature of outward normal n, I) - e0E + Pe (in SI units),
I) E + 4jiPe (in cgse units).
These formulas are generalizations of the formulas in Section \ssect{15.3.1}. Making use of the results obtained in Section \ssect{18.2.3}, we have lor an isotropic homogeneous medium
I) e0E+e0xeE=80(1+Xe)E(inSIunits), I) E + 4jxxeE = (1 + 4jiXe) E (in cgse units).
(consequently
I)—e0erE,whereer=1+Xe (inSIunits), I)—erE, where er.= 1+ 4jtxe (in cgse units).
The quantity er is the relative permittivity (dielectric con- slant) (cl*. \ssect{14.2.4}).
For vacuum er = 1 and Xe = 0.
I'he concept of *er introduced in Section \ssect{14.2.4} has a meaning only lor isotropic homogeneous media.

\section{Ferroelectric Materials}
\label{sec-18.4}


\subsection{}\label{18.4.1} Ferroelectrics belong to a group of crystalline dielectrics. I'he first substance of this type to be investigated was Seignette’s
salt (potassium sodium tartrate) NaKC4H406-4H20 and, con­ sequently, they are sometimes called seignette-electrics. An­ other example of a ferroelectric is barium titanate BaTi03. Typical of ferroelectrics is a drastic increase in the relative permittivity (dielectric constant) \ssect{14.2.4} in a definite tem­ perature range (Fig. 18.4).

\subsection{}\label{18.4.2} The relative permittivity er and the dielectric suscepti- bility Xe \ssect{18.2.4} of ferroelectrics are functions of the field strength
l 7—0186
E in the substance (Fig. 18.5). As a result, no linear relation between vectors Pe and E is observed in those substances. The dependence of the electric displacement I) on the held strength
Fig. 18.4 Fig. 18.5
is of a complex nature and a linear relation between D and E exists only at very high E values (Fig. 18.6).

\subsection{}\label{18.4.3} A monocrystal of a ferroelectric is divided into spontane­ ously polarized regions, called domains (cf. \ssect{26.5.4}). This spon­ taneous polarization of the domains is a result of the orientation of the dipole moments of the molecules in the domain in a defi-mlc direction. In the absence of an external electric field, the polarization vectors of the various domains are directed at ran­dom and for a large monocrystal or a polycrystal the total polarization is equal, on an average, to zero. When a ferro- electric is placed in an electric field, the electric moments of ilie domains are reoriented and the crystal acquires a total polarization not equal to zero.

\subsection{}\label{18.4.4} Domains are formed in ferroelectrics within a definite temperature range, between the upper and lower Curie points
()[‘p and 0jPw(cf. \ssect{26.5.2}). For Seignette’s salt (also called Rochelle salt) (18.4.1), 0qP= 298K and 0j2w= 258K. At T> 0gp
the forces of interaction between the dipoles cannot resist the ihermal motion, the spontaneous polarization of the domains is violated and the ferroelectric is converted into an ordinary polar dielectric. The drastic increase in heat capacity of the substance indicates that a second-order phase transition \ssect{12.4.2}) occurs at the Curie point. A disordered phase exists above the
(’.uric point 0qP and the dielectric is not polarized in the absence nl an external field. At T < 0£p there is an ordered phase indi­
cated by the spontaneous po­
larization in the domains*.

\subsection{}\label{18.4.5} The phenomenon of di­
electric, or ferromagnetic,hystere­
sis (lag) is observed in ferro-
olcctrics. As is evident from
Fig. 18.7, with an increase in the
strength of the external field,
the magnitude of vector Pe
increases and reaches satura­
tion at point a. Then, as E is
reduced to zero, a ferroelectric
lias remanent polarization, char­
acterized by the value of the polarization vector Pe0. Polarization completely disappears only after applying a depo­ larization electric field of the strength —Ec in the opposite
* A discussion of the phenomena that occur in ferroelectrics at T — 0!w and lead to the destruction of the domains is beyond the scope of the present handbook.
17*

direction. This is called the coercive force. Periodic reversal of polarization of a ferromagnetic material is accompanied by the consumption of electric energy, spent in heating the sub­ stance. The area of the hysteresis loop is proportional to the electric energy that is converted into internal energy per unit volume of the ferroelectric per cycle.